% ===================================================================
%  Metadatos: titulo, autores, abstract, keywords
%  Entregable: short-paper-es
%  Reestructurado 2026-09-06. Reglas en LATEX/CLAUDE.md
% ===================================================================
%% Autores y ORCID identicos en TODOS los entregables.
%% Fuente unica de verdad: CLAUDE.md, seccion 7 bis (Autoria).
%% El TITULO, ABSTRACT y KEYWORDS de abajo son heredados del ES: ajustarlos
%% a lo que pida este formato/revista.

\renewcommand{\figurename}{Fig.}
\renewcommand{\tablename}{TABLA}
\renewcommand{\refname}{REFERENCIAS}
% Secciones I..V, subsecciones A., B. y sub-subsecciones 1), 2), como el
% formato IEEE de conferencia. Sin esto, IEEEtran numera "III-B1.".
\renewcommand{\thesubsubsection}{\arabic{subsubsection}}
\renewcommand{\theparagraph}{\alph{paragraph}}

\title{Antes del Click: Análisis del Lado del Cliente para Mitigar la Probabilidad de Riesgo de Estafa en Compras Digitales}

\author{
\IEEEauthorblockN{Mathias Alejandro Jave Diaz}
\IEEEauthorblockA{\textit{Ingeniería de Software} \\
\textit{Universidad Peruana de} \\
\textit{Ciencias Aplicadas} \\
Lima, Perú \\
0009-0009-3705-505X}
\and
\IEEEauthorblockN{Alex Ramon Alberto Avila Asto}
\IEEEauthorblockA{\textit{Ingeniería de Software} \\
\textit{Universidad Peruana de} \\
\textit{Ciencias Aplicadas} \\
Lima, Perú \\
0009-0008-5624-8482}
\and
\IEEEauthorblockN{Fidel Eugenio Garcia Rojas}
\IEEEauthorblockA{\textit{Ingeniería de Software} \\
\textit{Universidad Peruana de} \\
\textit{Ciencias Aplicadas} \\
Lima, Perú \\
0009-0007-4364-0058}
}

\maketitle

\begin{abstract}
Los métodos eficaces contra el fraude en el comercio electrónico operan del lado de la plataforma utilizando telemetría interna, dejando al consumidor desprotegido en el momento de decisión. La evidencia indica que la autoconciencia no evita la victimización y que las verificaciones de legitimidad suelen realizarse después de efectuar el pago. Para abordar esta brecha, este artículo propone una arquitectura basada en una extensión de navegador que calcula un nivel estimado de riesgo de estafa en el equipo del consumidor, entendido como una escala ordinal y no como una probabilidad calibrada. El sistema resuelve ocho características (como la veracidad del dominio y la autenticidad de reseñas) a partir del documento HTML ya renderizado, sin enviar peticiones adicionales a la plataforma del vendedor, y tres de ellas se derivan en el motor de riesgo. Se incorpora una regla de decisión que distingue el riesgo detectado de la evidencia insuficiente, asegurando que la falta de información en una ficha de producto nunca se presente como seguridad. La evaluación diagnóstica sobre el conjunto de datos disponible reporta que cinco algoritmos de familias distintas clasifican la partición de prueba sin error: el conjunto es separable por construcción, de modo que el ejercicio vale como verificación del canal de procesamiento y no como medida de eficacia. El estudio define además el diseño experimental con el que la herramienta se someterá a prueba. Lo aportado es, por tanto, una propuesta arquitectónica con validación diagnóstica preliminar y protocolo experimental, y no evidencia de que el sistema mitigue la probabilidad de riesgo a la que se expone un comprador real.
\end{abstract}

\begin{IEEEkeywords}
comercio electrónico, compras digitales, extensión de navegador, nivel estimado de riesgo, riesgo del lado del cliente
\end{IEEEkeywords}
